\documentclass[10pt,a4paper]{amsart}
\usepackage[margin=19mm]{geometry}
\usepackage{amsmath,amssymb,mathtools}
\usepackage[scr=rsfs,cal=euler]{mathalfa}
\usepackage{xfrac}
\usepackage[unicode,hidelinks,bookmarks=false]{hyperref}
\newcommand{\ie}{i.e.\ }
\newcommand{\cf}{cf.\ }
\newcommand{\resp}{respectively\ }
\newcommand{\Subsection}{\subsection}
\providecommand{\nobreakdash}{\nobreak\hbox{-}}
\setlength{\emergencystretch}{2em}
\multlinegap0pt
\usepackage{amssymb}
\newtheorem{thm}{Theorem}
\newtheorem{prop}{Proposition}
\title{Special Solutions of $q$-Heun Equation by $q$-Integral Transformations}
\author{Ayaka Murakami, Kouichi Takemura}
\date{Source: arXiv:2605.01406v2 (CC BY 4.0)}
\begin{document}
\maketitle

\noindent\emph{Source and licence.} Special Solutions of $q$-Heun Equation by $q$-Integral Transformations. Version arXiv:2605.01406v2 (CC BY 4.0), distributed by arXiv under the Creative Commons Attribution 4.0 International licence, http://creativecommons.org/licenses/by/4.0/, as recorded in the arXiv licence metadata for this exact version and shown on the abstract page https://arxiv.org/abs/2605.01406v2. Full licence notice: \texttt{licenses/arxiv-2605.01406-CC-BY-4.0.txt}.\par\smallskip

\noindent\textit{Editorial scope.} Counted unit: the article's prop prop (source0.tex lines 1054--1057). The article's complete original proof of it is delivered with it (source0.tex lines 1058--1102, one \texttt{proof} environment of 45 lines). No mathematics is written, repaired or re-derived: the statement and the proof are the article's own lines, in the article's own notation, and no retained line is substituted or re-flowed.  Retained uncounted as the source-local context the proof needs: lines 409--411; lines 946--1035; lines 951--955.  Nothing was omitted from any retained span, so the package carries no omission record.  No retained line contains a citation, so the wrapper carries no bibliography and no bibliography entry.  No retained line contains a figure environment, an image-inclusion command or drawing code, so no image is delivered and the package declares no asset dependency.  Editorial formatting only, in addition: an AMS \texttt{amsart} preamble that restates the article's own declarations and macros used by the retained lines, namely the environments \texttt{thm}, \texttt{prop} on separate counters, together with the standard packages those lines need.  The retained environments print in this document as Theorem 1, Proposition 1 in order of appearance, whereas the article prints the same items with the numbering of its own full document.  The complete native source of the article as distributed in the arXiv e-print for this exact version is bundled unchanged as \texttt{sources/199634/source0.tex}.
\begin{align}
 & c (E'') = \sum (-1) ^k x''_{i_1} z'' _{i_1 +1 } x''_{i_2} z'' _{i_2 +1 } \cdots x''_{i_k} z'' _{i_k +1 }  \prod _{j \in \{ 1,2, \dots ,N +1 \} \setminus \atop{\{ i_1 , i_1 +1 , \dots , i_k , i_k +1  \} }} (E'' + y''_j ) , \label{eq:cE''sum}
\end{align}

\begin{thm} \label{thm:2sols}
%Assume that $\beta = N +1 $ and $N$ is a non-negative integer.
Set $\lambda_{1} = (h_{1} + h_{2} - l_{1} - l_{2} - \alpha_{1} - \alpha_{2} - \beta + 2)/2$.
Let $N$ be a non-negative integer, and assume that $\beta = N +1 $ and $-\lambda_{1} -\alpha _1 \not \in \{ 0,1,\dots ,N-1 \}$.
Set $c_{-1}(E)=0 $,  $c_0(E)=1 $,
\begin{align}
 & x_n =t_{1}t_{2} q^{ n - N +1 + h_{1} + h_{2} - \alpha _1 - 2 \lambda _1 } (1 - q^{-n})(1 - q^{N-n + \lambda _1 + \alpha _1 }) , \label{eq:2xyz}  \\
 & y_n =  q^{N -n + 1/2 + \lambda _1 + \alpha_{1} + \alpha_{2} } (q^{l_{1}}t_{1} + q^{l_{2}}t_{2}) + q^{n- N-1/2  - \lambda _1} (q^{h_{1}}t_{1} + q^{h_{2}}t_{2})  , \notag \\
 & z_n = q^{ n - N-2 + \alpha_{1}  } (1 - q^{N -n +1 +\lambda _1 + \alpha_{2}  })(1 - q^{N- n + 2}), \notag
\end{align}
and we determine the polynomials $c_n(E)$ $(n=1,\dots ,N)$ recursively by
\begin{align}
 & c_{n} (E) x_n = c_{n - 1} (E) [E + y_n ] -c_{n-2} (E) z_n . \label{eq:2cnE}
\end{align}
We define the polynomial $c (E)$ by
\begin{align}
 & c (E) = x_1 x_2 \cdots x_{N} [ c_{N} (E) \{ E + y_{N+1} \} -c_{N-1} (E) z_{N+1} ] . \label{eq:2cE}
\end{align}
Assume that $E=E_0$ is a solution of the algebraic equation $c (E)=0$.
\\
(i) Let $\xi $ be the variable which is independent of the variable $x$ or proportional to $x$.
% (i.~e.~$\xi =Ax$ where $A $ is independent of $x$).
%Then, 
If  $ \lambda _{1} + \alpha _{2}  > 1$, then the Jackson integral
  \begin{align*}
    &g _1 (x) = (1 - q) x^{-\alpha_{1}} \frac{(q^{-l_{1} + 1/2}\xi/t_{1}, q^{-l_{2} + 1/2}\xi/t_{2}, q^{ \lambda _{1} + \alpha _{1} }\xi/x; q)_{\infty}}{(q^{ \lambda _{1} -h_{1} + \alpha _{1} - 1/2}\xi/t_{1}, q^{ \lambda _{1} -h_{2}+ \alpha _{1} - 1/2}\xi/t_{2}, \xi/x; q)_{\infty}} \\
    &\qquad \times \sum_{n = -\infty}^{\infty} \frac{(q^{ \lambda _{1} -h_{1} + \alpha _{1} - 1/2}\xi/t_{1}, q^{ \lambda _{1} -h_{2} + \alpha _{1} - 1/2}\xi/t_{2}, \xi/x; q)_{n}}{(q^{-l_{1} + 1/2}\xi/t_{1}, q^{-l_{2} + 1/2}\xi/t_{2}, q^{ \lambda _{1} + \alpha _{1} }\xi/x; q)_{n}} \sum_{k = 0}^{N} q^{( k+1)n} \xi^{k+1} c_{k}(E_0) \notag
  \end{align*}
converges and it satisfies
  \begin{align}
    &A^{\langle 4 \rangle} (x; h_{1}, h_{2}, l_{1}, l_{2}, \alpha_{1}, \alpha_{2}, \beta) g (x) \label{eq:2qHnonh1} \\
    &\qquad = E_0 g (x) - (1 - q) x^{- \alpha_{1}}q^{ - \lambda _{1} + h_{1} + h_{2} + 1 }(q^{- \lambda _{1} - \alpha _{1} - N} - 1)t_{1}t_{2} . \notag
  \end{align}
(ii) If  $ \lambda _{1} + \alpha _{2}  > 1$, then the Jackson integral
  \begin{align*}
    &g _2 (x) = (1 - q) x^{ \lambda _{1} } \frac{(q^{-  \lambda _{1} + h_{1} - \alpha _{1} + 3/2}t_{1}/\xi, q^{ - \lambda _{1} +h_{2} - \alpha _{1} + 3/2}t_{2}/\xi, qx/\xi; q)_{\infty}}{(q^{l_{1} + 1/2}t_{1}/\xi, q^{l_{2} + 1/2}t_{2}/\xi, q^{- \lambda _{1} - \alpha _{1} + 1}x/\xi; q)_{\infty}} \\
    & \qquad \qquad \times \sum_{n = -\infty}^{\infty} \frac{(q^{l_{1} + 1/2}t_{1}/\xi, q^{l_{2} + 1/2}t_{2}/\xi, q^{- \lambda _{1} - \alpha _{1} + 1 }x/\xi; q)_{n}}{(q^{- \lambda _{1} + h_{1} - \alpha _{1} + 3/2}t_{1}/\xi, q^{- \lambda _{1} + h_{2} - \alpha _{1} + 3/2}t_{2}/\xi, qx/\xi; q)_{n}} \\
    & \qquad \qquad \qquad \qquad \times \sum_{k = 0}^{N} q^{( \lambda _{1}  + \alpha_{2}+ N -k)n} \xi^{- \lambda _{1}  -\alpha_{2} - N + k} c_{k}(E_0) \notag
%    &g _2 (x) = (1 - q)\xi^{- \lambda _{1}  -\alpha_{2} - N }x^{ \lambda _{1} } \frac{(q^{-  \lambda _{1} + h_{1} - \alpha _{1} + 3/2}t_{1}/\xi, q^{ - \lambda _{1} +h_{2} - \alpha _{1} + 3/2}t_{2}/\xi, qx/\xi; q)_{\infty}}{(q^{l_{1} + 1/2}t_{1}/\xi, q^{l_{2} + 1/2}t_{2}/\xi, q^{- \lambda _{1} - \alpha _{1} + 1}x/\xi; q)_{\infty}} \\
%    &\times \sum_{n = -\infty}^{\infty} \frac{(q^{l_{1} + 1/2}t_{1}/\xi, q^{l_{2} + 1/2}t_{2}/\xi, q^{- \lambda _{1} - \alpha _{1} + 1 }x/\xi; q)_{n}}{(q^{- \lambda _{1} + h_{1} - \alpha _{1} + 3/2}t_{1}/\xi, q^{- \lambda _{1} + h_{2} - \alpha _{1} + 3/2}t_{2}/\xi, qx/\xi; q)_{n}} q^{( \lambda _{1}  + \alpha_{2}+ N )n} \sum_{k = 0}^{N} q^{-nk} \xi^{k} c_{k}(E_0) \notag
  \end{align*}
converges and it satisfies
  \begin{align*}
    & A^{\langle 4 \rangle} (x; h_{1}, h_{2}, l_{1}, l_{2}, \alpha_{1}, \alpha_{2}, \beta) g(x) = E_0 g(x) - (1 - q) x^{ \lambda _{1} } q^{ - \lambda _{1} + h_{1} + h_{2} + 1} \xi^{- \lambda _{1} -\alpha_{2}  - N  -1}  \\
    & \quad \times \frac{\vartheta _q (q^{- \lambda _{1} + h_{1} - \alpha _{1} + 3/2}t_{1}/\xi) \vartheta _q (q^{- \lambda _{1} + h_{2} - \alpha _{1} + 3/2}t_{2}/\xi) \vartheta_{q} \big(qx/\xi \big)}{\vartheta _q (q^{l_{1} + 1/2}t_{1}/\xi) \vartheta _q (q^{l_{2} + 1/2}t_{2}/\xi) \vartheta_{q} (q^{- \lambda _{1} - \alpha _{1} + 1 } x/\xi )} \big( q^{ - \lambda _{1} - \alpha _{1} -N} - 1 \big) t_{1} t_{2} . \notag
  \end{align*}
(iii) Set
\begin{align*}
& g_{3}(x) = x^{ \lambda _{1} } \frac{( q^{ \lambda _{1} - h_{1} + \alpha _{1} + 1/2}x/t_{1}; q)_\infty}{( q^{-h_{1} + 1/2}x/t_{1}; q)_\infty}  \sum_{k = 0}^{N} (q^{- \lambda _{1} + h_{1}  - \alpha _{1} + 1/2}t_{1})^k c_{k} (E_{0})  \notag \\
  &\quad \times {}_3\phi_2 \biggl( \begin{array}{@{}c@{}} q^{ \lambda _{1} -h_{1} + l_{1} + \alpha _{1} }, q^{ \lambda _{1} -h_{1} + l_{2} + \alpha _{1} }t_{2}/t_{1}, q^{-h_{1} + 1/2}x/t_{1} \\
     q^{- h_{1} + h_{2} + 1}t_{2}/t_{1}, q^{ \lambda _{1} - h_{1} + \alpha _{1} + 1/2}x/t_{1} \end{array} ;q,q^{ \lambda _{1}  + \alpha_{2} +N  - k} \biggr), \\
  &  g_{4}(x) = x^{ \lambda _{1} } \frac{( q^{ \lambda _{1} - h_{2} + \alpha _{1} + 1/2}x/t_{2}; q)_\infty}{( q^{-h_{2} + 1/2}x/t_{2}; q)_\infty}  \sum_{k = 0}^{N}  (q^{- \lambda _{1} +h_{2} - \alpha _{1} + 1/2}t_{2})^k c_{k} (E_{0}) \notag \\
  &\quad \times {}_3\phi_2 \biggl( \begin{array}{@{}c@{}} q^{ \lambda _{1} -h_{2} + l_{1} + \alpha _{1} }t_{1}/t_{2}, q^{ \lambda _{1} -h_{2} + l_{2} + \alpha _{1} }, q^{-h_{2} + 1/2}x/t_{2} \\
     q^{h_{1} - h_{2} + 1}t_{1}/t_{2}, q^{ \lambda _{1} - h_{2} + \alpha _{1} + 1/2}x/t_{2} \end{array} ;q,q^{ \lambda _{1}  + \alpha_{2} +N - k} \biggr), \\
 & g_{5}(x) = x^{ - \alpha_{2} } \frac{(q^{- \lambda _{1} + h_{1} - \alpha _{1} + 3/2}t_{1}/x, q^{- \lambda _{1} + h_{2} - \alpha _{1} + 3/2}t_{2}/x ; q)_{\infty}}{(q^{l_{1} + 1/2}t_{1}/x, q^{l_{2} + 1/2}t_{2}/x ; q)_{\infty}} \sum_{k = 0}^{N} x^{k-N} c_{k}  (E_{0}) \notag \\
  &\quad \times  {}_3\phi_2 \biggl( \begin{array}{@{}c@{}} q^{l_{1} + 1/2}t_{1}/x, q^{l_{2} + 1/2}t_{2}/x, q^{- \lambda _{1} - \alpha _{1} + 1} \\
     q^{- \lambda _{1} + h_{1} - \alpha _{1} + 3/2}t_{1}/x, q^{- \lambda _{1} + h_{2} - \alpha _{1} + 3/2}t_{2}/x \end{array} ;q,q^{ \lambda _{1}  + \alpha_{2} +N - k} \biggr).
%& g_{3}(x) = (1 - q) x^{ \lambda _{1} } \frac{(q, q^{- h_{1} + h_{2} + 1}t_{2}/t_{1}, q^{ \lambda _{1} - h_{1} + \alpha _{1} + 1/2}x/t_{1}; q)_\infty}{(q^{ \lambda _{1} -h_{1} + l_{1} + \alpha _{1} }, q^{ \lambda _{1} -h_{1} + l_{2} + \alpha _{1} }t_{2}/t_{1}, q^{-h_{1} + 1/2}x/t_{1}; q)_\infty} \notag \\
%  &\quad \times \sum_{k = 0}^{N} c_{k} (E_{0}) (q^{- \lambda _{1} + h_{1}  - \alpha _{1} + 1/2}t_{1})^{- \lambda _{1} - \alpha_{2} - N + k} \notag \\
%  &\quad \times {}_3\phi_2 \biggl( \begin{array}{@{}c@{}} q^{ \lambda _{1} -h_{1} + l_{1} + \alpha _{1} }, q^{ \lambda _{1} -h_{1} + l_{2} + \alpha _{1} }t_{2}/t_{1}, q^{-h_{1} + 1/2}x/t_{1} \\
%     q^{- h_{1} + h_{2} + 1}t_{2}/t_{1}, q^{ \lambda _{1} - h_{1} + \alpha _{1} + 1/2}x/t_{1} \end{array} ;q,q^{ \lambda _{1}  + \alpha_{2} +N  - k} \biggr), \\
%  &  g_{4}(x) = (1 - q) x^{ \lambda _{1} } \frac{(q^{h_{1} - h_{2} + 1}t_{1}/t_{2}, q, q^{ \lambda _{1} - h_{2} + \alpha _{1} + 1/2}x/t_{2}; q)_\infty}{(q^{ \lambda _{1} -h_{2} + l_{1} + \alpha _{1} }t_{1}/t_{2}, q^{ \lambda _{1} -h_{2} + l_{2} + \alpha _{1} }, q^{-h_{2} + 1/2}x/t_{2}; q)_\infty} \notag \\
%  &\quad \times \sum_{k = 0}^{N} c_{k} (E_{0}) (q^{- \lambda _{1} +h_{2} - \alpha _{1} + 1/2}t_{2})^{ -\lambda _{1} - \alpha_{2} -N  + k} \notag \\
%  &\quad \times {}_3\phi_2 \biggl( \begin{array}{@{}c@{}} q^{ \lambda _{1} -h_{2} + l_{1} + \alpha _{1} }t_{1}/t_{2}, q^{ \lambda _{1} -h_{2} + l_{2} + \alpha _{1} }, q^{-h_{2} + 1/2}x/t_{2} \\
%     q^{h_{1} - h_{2} + 1}t_{1}/t_{2}, q^{ \lambda _{1} - h_{2} + \alpha _{1} + 1/2}x/t_{2} \end{array} ;q,q^{ \lambda _{1}  + \alpha_{2} +N - k} \biggr), \\
% & g_{5}(x) =(1 - q) x^{ - \alpha_{2} - N  } \frac{(q^{- \lambda _{1} + h_{1} - \alpha _{1} + 3/2}t_{1}/x, q^{- \lambda _{1} + h_{2} - \alpha _{1} + 3/2}t_{2}/x, q; q)_{\infty}}{(q^{l_{1} + 1/2}t_{1}/x, q^{l_{2} + 1/2}t_{2}/x, q^{- \lambda _{1} - \alpha _{1} + 1}; q)_{\infty}} \notag \\
%  &\quad \times \sum_{k = 0}^{N} x^k c_{k}  (E_{0}) {}_3\phi_2 \biggl( \begin{array}{@{}c@{}} q^{l_{1} + 1/2}t_{1}/x, q^{l_{2} + 1/2}t_{2}/x, q^{- \lambda _{1} - \alpha _{1} + 1} \\
%     q^{- \lambda _{1} + h_{1} - \alpha _{1} + 3/2}t_{1}/x, q^{- \lambda _{1} + h_{2} - \alpha _{1} + 3/2}t_{2}/x \end{array} ;q,q^{ \lambda _{1}  + \alpha_{2} +N - k} \biggr).  
\end{align*}
The functions $g_3(x)$,  $g_4(x)$ and $g_5(x)$ satisfy the $q$-Heun equation
  \begin{align*}
    & A^{\langle 4 \rangle} (x; h_{1}, h_{2}, l_{1}, l_{2}, \alpha_{1}, \alpha_{2}, \beta) g(x) = E_0 g(x) . 
  \end{align*}
(iv)
Set 
  \begin{align*}
    &g_{6}(x) = (1 - q)x^{-\alpha_{1}}\frac{(q, q^{l_{1} - l_{2} + 1}t_{1}/t_{2}, q^{ \lambda _{1} + l_{1} +\alpha _{1} + 1/2}t_{1}/x; q)_{\infty}}{(q^{ \lambda _{1} -h_{1} + l_{1} + \alpha _{1} }, q^{ \lambda _{1} -h_{2} + l_{1} + \alpha _{1} }t_{1}/t_{2}, q^{l_{1} + 1/2}t_{1}/x; q)_{\infty}} \\
    &\quad \times \sum_{k = 0}^{N} (q^{l_{1} + 1/2}t_{1})^{k+1} c_k (E_{0}) {}_3\phi_2 \biggl( \begin{array}{@{}c@{}} q^{ \lambda _{1} -h_{1} + l_{1} + \alpha _{1} }, q^{ \lambda _{1} -h_{2} + l_{1} + \alpha _{1} }t_{1}/t_{2}, q^{l_{1} + 1/2}t_{1}/x \\
q^{l_{1} - l_{2} + 1}t_{1}/t_{2}, q^{\lambda _{1} + l_{1} + \alpha _{1} + 1/2}t_{1}/x \end{array} ;q,q^{k+1} \biggr), \\
    &g_{7}(x) = (1 - q)x^{-\alpha_{1}}\frac{(q,q^{- l_{1} + l_{2} + 1}t_{2}/t_{1}, q^{\lambda _{1} + l_{2} + \alpha _{1} + 1/2}t_{2}/x; q)_{\infty}}{(q^{ \lambda _{1} -h_{2} + l_{2} + \alpha _{1} }, q^{ \lambda _{1} -h_{1} + l_{2} + \alpha _{1} }t_{2}/t_{1},  q^{l_{2} + 1/2}t_{2}/x; q)_{\infty}} \\
    &\quad \times \sum_{k = 0}^{N} (q^{l_{2} + 1/2}t_{2})^{k+1}  c_k (E_{0}) {}_3\phi_2 \biggl( \begin{array}{@{}c@{}} q^{ \lambda _{1} -h_{2} + l_{2} + \alpha _{1} }, q^{ \lambda _{1} -h_{1} + l_{2}  + \alpha _{1} }t_{2}/t_{1},  q^{l_{2} + 1/2}t_{2}/x \\
q^{- l_{1} + l_{2} + 1}t_{2}/t_{1}, q^{\lambda _{1} + l_{2} + \alpha _{1} + 1/2}t_{2}/x \end{array} ;q,q^{k+1} \biggr), \\
    &g_{8}(x) = (1 - q)x^{-\alpha_{1}}\frac{(q, q^{- \lambda _{1} -l_{1} - \alpha _{1} + 3/2}x/t_{1}, q^{- \lambda _{1} -l_{2}  - \alpha _{1} + 3/2}x/t_{2} ; q)_{\infty}}{( q^{- \lambda _{1} - \alpha _{1} + 1}, q^{-h_{1} + 1/2}x/t_{1}, q^{-h_{2} + 1/2}x/t_{2}; q)_{\infty}} \\
    &\quad \times \sum_{k = 0}^{N} (q^{- \lambda _{1} - \alpha _{1} + 1}x)^{k+1}  c_k (E_{0}) {}_3\phi_2 \biggl( \begin{array}{@{}c@{}} q^{- \lambda _{1} - \alpha _{1} + 1}, q^{-h_{1} + 1/2}x/t_{1}, q^{-h_{2} + 1/2}x/t_{2} \\
q^{- \lambda _{1} -l_{1} - \alpha _{1} + 3/2}x/t_{1}, q^{- \lambda _{1} -l_{2} - \alpha _{1} + 3/2}x/t_{2} \end{array} ;q,q^{k+1} \biggr).
  \end{align*}
The functions $g_{i} (x)- g_{j} (x)$ $(i, \, j \in \{6, 7, 8\})$ satisfy the $q$-Heun equation.
\\
(v) The functions $g_6(x)$,  $g_7(x)$ and $g_8(x)$ satisfy equation (\ref{eq:2qHnonh1}).
 \end{thm}

\begin{align}
 & x_n =t_{1}t_{2} q^{ n - N +1 + h_{1} + h_{2} - \alpha _1 - 2 \lambda _1 } (1 - q^{-n})(1 - q^{N-n + \lambda _1 + \alpha _1 }) ,   \\
 & y_n =  q^{N -n + 1/2 + \lambda _1 + \alpha_{1} + \alpha_{2} } (q^{l_{1}}t_{1} + q^{l_{2}}t_{2}) + q^{n- N-1/2  - \lambda _1} (q^{h_{1}}t_{1} + q^{h_{2}}t_{2})  , \notag \\
 & z_n = q^{ n - N-2 + \alpha_{1}  } (1 - q^{N -n +1 +\lambda _1 + \alpha_{2}  })(1 - q^{N- n + 2}), \notag
\end{align}

\begin{prop} 
Assume that $\beta = N +1 $ and $N$ is a non-negative integer.
The condition that the singularity $x=0$ of the $q$-Heun equation is apparent is equivalent to the condition that the eigenvalue $E$ satisfies the algebraic equation $c(E)=0$ in Theorem \ref{thm:2sols}.
\end{prop}

\begin{proof}
We recall the definition that the singularity $x=0$ is apparent in the setting of the $q$-Heun equation.
Set $\lambda_{1} = (h_{1} + h_{2} - l_{1} - l_{2} - \alpha_{1} - \alpha_{2} - \beta + 2)/2$.
If the formal series $ s^{\lambda _{1}} \sum_{n =0}^{\infty } d_{n} x^{n} $ is a solution to the $q$-Heun equation, then we have
\begin{align}
 & d_{n} t_{1}t_{2}q^{  1- n + h_{1} + h_{2} - \lambda _{1}} (1 - q^{n})(1 - q^{n - \beta }) \label{eq:recreldn} \\
 & \ - d_{n - 1} [E + q^{3/2- n - \lambda _{1} } (q^{h_{1}}t_{1} + q^{h_{2}}t_{2} ) + q^{n - 3/2 + \lambda _{1} + \alpha _{1} + \alpha _{2} }( q^{l_{1}}t_{1} + q^{l_{2}}t_{2} ) ] \notag \\
% & \ - d_{n - 1} [E'' + q^{3/2- n - \lambda ''_{1} +h''_{1}}t_{1} + q^{3/2- n - \lambda ''_{1} + h''_{2}}t_{2} + q^{n - 3/2 + \lambda ''_{1} + \alpha ''_{1} + \alpha ''_{2} + l''_{1}}t_{1} + q^{n - 3/2 + \lambda ''_{1} + \alpha ''_{1} + \alpha ''_{2} + l''_{2}}t_{2} ] \notag \\
 & \ + d_{n - 2} q^{ 2- n - \lambda _{1} } (1 - q^{n - 2 + \lambda _{1} + \alpha _{1} })(1 - q^{n - 2 + \lambda _{1} + \alpha _{2} }) = 0 . \notag
  \end{align}
for $n = 1 ,2,3, \cdots $, where $d_{-1} =0$.
Set $d_0=1$ and determine $d_n$ for $n=1,2,\dots ,N$ recursively by equation (\ref{eq:recreldn}).
On equation (\ref{eq:recreldn}) for $n=N+1$, the coefficint of $d_{N+1}$ vanishes, because $\beta = N +1 $.
Then, the apparency of the singularity $x=0$ is that the equation (\ref{eq:recreldn}) for $n=N+1$ holds.

Set $d_{-1}(E)=0 $,  $d_0(E)=1 $,
\begin{align*}
 & x^* _n = t_{1}t_{2}q^{  1- n + h_{1} + h_{2} - \lambda _{1}} (1 - q^{n})(1 - q^{n - \beta }) , \\
 & y^* _n = q^{3/2- n - \lambda _{1} } (q^{h_{1}}t_{1} + q^{h_{2}}t_{2} ) + q^{n - 3/2 + \lambda _{1} + \alpha _{1} + \alpha _{2} }( q^{l_{1}}t_{1} + q^{l_{2}}t_{2} ) , \notag \\
 & z^* _n = q^{ 2- n - \lambda _{1} } (1 - q^{n - 2 + \lambda _{1} + \alpha _{1} })(1 - q^{n - 2 + \lambda _{1} + \alpha _{2} }) , \notag
\end{align*}
and we determine the polynomials $d_n(E)$ $(n=1,\dots ,N)$ recursively by
\begin{align}
 & d_{n} (E) x^* _n = d_{n - 1} (E) [E + y^* _n ] -d_{n-2} (E) z^* _n . \label{eq:dnE}
\end{align}
We define the polynomial $d (E)$ by
\begin{align}
 & d (E) = x^* _1 x^* _2 \cdots x^* _{N} [ d_{N} (E) \{ E + y^* _{N+1} \} - d_{N-1} (E) z^* _{N+1} ] . \label{eq:dE}
\end{align}
Then, the singularity $x=0$ is apparent, if and only if the eigenvalue $E$ satisfies $d(E)=0$.

It is enough to show that $c(E) = d(E)$.
Let $x_n , y_n, z_n$ be the values defined in equation (\ref{eq:2xyz}).
Then, we have
\begin{equation*}
y^* _{N+2-n} =y_n , \; x^* _{N+1-n} z^* _{N+2-n} = x_n z_{n+1} .
\end{equation*}
It follows from equation (\ref{eq:cE''sum}) and the corresponding equation for $d(E)$ that
\begin{align*}
 & d (E) = \sum (-1) ^k x^* _{j_1} z^* _{j_1 +1 } x^* _{j_2} z^* _{j_2 +1 } \cdots x^* _{j_k} z^* _{j_k +1 } \! \! \! \! \! \! \! \! \! \!  \prod _{j' \in \{ 1,2, \dots ,N +1 \} \setminus \atop{\{ j_1 , j_1 +1 , \dots , j_k , j_k +1  \} }} \! \! \! \! \! \! \! \! \! \! (E + y^* _{j'} ) , \\
%\label{eq:dEsum}
& = \sum (-1) ^k x^* _{N+1 - i_k} z^* _{N +2 - i_k } \cdots x^* _{N+1 -i_2} z^* _{N+2 - i_2 } x^* _{N+1 -i_1} z^* _{N+2 -i_1 } \! \! \! \! \! \! \! \! \! \! \prod _{i' \in \{ 1,2, \dots ,N +1 \} \setminus \atop{\{ i_1 , i_1 +1 , \dots , i_k , i_k +1  \} }}  \! \! \! \! \! \! \! \! \! \!  (E + y^* _{N+2-i'} ) \notag \\
 &  = \sum (-1) ^k x _{i_1} z _{i_1 +1 } x _{i_2} z _{i_2 +1 } \cdots x _{i_k} z _{i_k +1 } \! \! \! \! \! \! \! \! \! \!  \prod _{i' \in \{ 1,2, \dots ,N +1 \} \setminus \atop{\{ i_1 , i_1 +1 , \dots , i_k , i_k +1  \} }} \! \! \! \! \! \! \! \! \! \! (E + y _{i'} ) = c(E) . \notag 
\end{align*}
\end{proof}

\end{document}
